\documentclass{article}
\usepackage[pdftex]{graphicx}
\usepackage{dot2texi}
\usepackage{tikz}
\usetikzlibrary{shapes,arrows}

\begin{document}

\section*{Factor Graph Representation of Trivial SLAM}
This graph represents a 3-pose, 3-landmark SLAM problem. 
Circles represent \textbf{Variable Nodes} (State) and Squares represent \textbf{Factor Nodes} (Constraints).

\begin{dot2tex}[dot,options=-lneato]
graph G {
    rankdir=LR;
    nodesep=0.5;
    
    // Variable Nodes: Poses (Circles)
    node [shape=circle, style=filled, fillcolor=lightblue];
    PoseA [label="Pose A"];
    PoseB [label="Pose B"];
    PoseC [label="Pose C"];

    // Variable Nodes: Landmarks (Circles)
    node [shape=circle, style=filled, fillcolor=lightyellow];
    LandmarkMug [label="Mug"];
    LandmarkTV [label="TV"];
    LandmarkLamp [label="Lamp"];

    // Factor Nodes: Motion (Squares)
    node [shape=square, style=filled, fillcolor=salmon, fixedsize=true, width=0.2, label=""];
    
    // Factors connecting Poses (Odometry)
    PriorFactor [xlabel="Prior"];
    OdomAB [xlabel="Odom AB"];
    OdomBC [xlabel="Odom BC"];

    // Wiring Motion
    PriorFactor -- PoseA;
    PoseA -- OdomAB -- PoseB;
    PoseB -- OdomBC -- PoseC;

    // Factor Nodes: Vision (Small Squares)
    node [shape=square, style=filled, fillcolor=lightgreen, width=0.15];

    // Pose A observations
    FactAMug -- PoseA; FactAMug -- LandmarkMug;
    FactATV -- PoseA; FactATV -- LandmarkTV;

    // Pose B observations
    FactBTV -- PoseB; FactBTV -- LandmarkTV;
    FactBLamp -- PoseB; FactBLamp -- LandmarkLamp;

    // Pose C observations
    FactCLamp -- PoseC; FactCLamp -- LandmarkLamp;
    
    // The Loop Closure: Pose C seeing the first landmark again
    FactCMug [fillcolor=red, xlabel="Loop Closure"];
    FactCMug -- PoseC; 
    FactCMug -- LandmarkMug [color=red, penwidth=2.0];
}
\end{dot2tex}

\end{document}
